\section{Reproducing the Paper's Attack and Testing the Repair}\label{sec:repro}

\subsection{The malicious-KGC forgery on Verma's scheme}\label{sec:repro-attack}
\prov{Paper}, reproduced. The attacker is the paper's $\Ftwo$, a malicious KGC. It issued sensor $i$'s
certificate $(R_{1i},c_i)$ and has observed one valid signature $(R_i,z_i)$ on a message $m_i$. It does
not know $sk_i$ or the nonce $r_{2i}$, and it wants a signature on a new message $m'$.

From \cref{eq:verma-sign}, the KGC can strip its own contributions: $R_i-R_{1i}=r_{2i}P$, and
$z_i-c_i=r_{2i}+sk_iv_i$, a single scalar that mixes the two unknowns. The weakness is that the hash
$v_i=H_1(m_i\concat pk_i\concat id_i\concat\Delta)$ does not depend on the nonce, so for any new message
the coefficient $v'=H_1(m'\concat pk_i\concat id_i\concat\Delta)$ can be computed in advance. Rescaling by
$\lambda=v'/v_i$ gives
\begin{equation}\label{eq:rescale}
R'=R_{1i}+\lambda\,(R_i-R_{1i}),\qquad z'=c_i+\lambda\,(z_i-c_i),
\end{equation}
which is a valid signature on $m'$ with the unknown nonce $\lambda r_{2i}$: substituting into
\cref{eq:verma-verify} gives $z'P=R'+h_0^i\PTA+v'pk_i$. The attacker never learns $r_{2i}$ or $sk_i$; it
rescales a value it cannot open. Our function \code{forge\_verma} follows the paper's derivation step by
step and takes no secret key as an argument, which a test checks.

\subsection{Results}
We ran the attack for 4 target messages on each backend, against Verma's scheme and, with the same
strategy, against the repaired scheme (\cref{tab:repro}). Against Verma's scheme every forgery verifies,
also inside an aggregate next to an honest signature. Against the repair the rescaling has no fixed
target: the new commitment $T'$ must be chosen before $v'$ can be computed, but the attack defines $T'$
in terms of $v'$. Iterating this circular definition for 16 rounds produced 16 distinct candidates in
every trial and no forgery.

\begin{table}[ht]
\centering\small
\caption{The paper's attack and our ablation (4 messages $\times$ 2 backends).}\label{tab:repro}
\begin{tabular}{llcl}
\toprule
& Scheme & Forgeries that verify & Status \\
\midrule
A & Verma \emph{et al.} & 8 / 8 & broken, as the paper says \prov{Reproduced} \\
B & Qiao \emph{et al.}, repaired & 0 / 8 & the attack does not apply \prov{Reproduced} \\
C & Qiao without the $T_i$ binding & 8 / 8 & broken again \prov{Ours} \\
\bottomrule
\end{tabular}
\end{table}

\subsection{Which change stops the attack}
The repair changes several things at once (\cref{sec:schemes}), so its resistance alone does not say
which change matters. We built a control: the repaired scheme with exactly one change removed, the
nonce commitment $T_i$ dropped from $H_1$ and $H_2$. Everything else is the real scheme's code: the same
\textsf{Setup}, \textsf{KeyGen}, \textsf{CertGen}, aggregation and verification equation, with $R_i$ still
hashed and $u_i,v_i$ still from separate hash functions. The variant is a working scheme: honest
signatures verify and signatures do not transfer to other messages (8/8). The paper's attack succeeds
against it in 8 of 8 trials (row C).

Binding $T_i$ into the hash is therefore \emph{necessary} to stop this attack, all else held fixed. The
paper states the same design rule in its Section~V-A: the commitment $T_i$ ``should be used to generate
the core element of signature'', otherwise a malicious KGC can forge from a known
signature~\cite{qiao2023}. The ablation tests that rule; it does not show that the other changes are
unnecessary, nor that the repaired scheme is secure, which is the paper's proof. \Cref{sec:keyonly}
returns to the other binding.
